\chapter{Resumo}
\label{CH:SUM}

Este texto apresentou as principais ideias em sistemas operacionais por meio do estudo de um sistema operacional, o xv6, linha por linha. Algumas linhas de código incorporam a essência das ideias 
principais (por exemplo, troca de contexto, fronteira entre usuário e kernel, travas, etc.) e cada linha é importante; outras linhas de código servem como ilustração de como implementar uma determinada 
ideia de sistema operacional e poderiam ser feitas de maneiras diferentes (por exemplo, um algoritmo melhor para escalonamento, estruturas de dados melhores no disco para representar arquivos, melhor 
logging para permitir transações concorrentes, etc.). Todas as ideias foram ilustradas no contexto de uma interface de chamadas de sistema específica e muito bem-sucedida, a interface Unix, mas essas 
ideias se aplicam ao projeto de outros sistemas operacionais.
